\chapter{Magnesium deficiency}

\section*{Definition and Overview}
Magnesium deficiency, clinically termed hypomagnesaemia, is an electrolyte disorder characterised by a subnormal concentration of magnesium in the blood (typically defined as a serum magnesium level below 0.7 mmol/L). Magnesium is the second most abundant intracellular cation and is an essential cofactor in over 300 enzymatic reactions. It plays a critical role in cellular energy production (ATP synthesis), DNA and RNA replication, protein synthesis, muscle contraction, nerve conduction, blood pressure regulation, and cardiac rhythm stability. Because less than 1\% of the body's total magnesium stores is present in the extracellular fluid, a patient can have a significant cellular magnesium depletion despite presenting with a normal serum magnesium level, a condition known as subclinical magnesium deficiency.

\section*{Epidemiology}
Hypomagnesaemia is a common clinical finding, particularly in hospitalised and critically ill populations. It is estimated to affect 10--20\% of hospitalised ward patients and up to 50--60\% of patients in intensive care units (ICUs). In the general community, subclinical magnesium deficiency is highly prevalent, often due to modern dietary patterns dominated by processed foods and a low intake of green vegetables and whole grains. High-risk groups include chronic alcoholics, elderly individuals, patients with type 2 diabetes mellitus, and those with chronic gastrointestinal disorders. There is no significant gender or racial disparity in its incidence, though it correlates strongly with comorbid chronic illnesses and polypharmacy.

\section*{Aetiology and Pathophysiology}
Magnesium deficiency arises from inadequate intake, impaired gastrointestinal absorption, or accelerated renal excretion. The primary causes and pathophysiology include:
\begin{itemize}
    \item \textbf{Inadequate dietary intake:} A diet poor in unprocessed foods, whole grains, leafy green vegetables, and nuts, which are the primary dietary sources of magnesium.
    \item \textbf{Gastrointestinal malabsorption:} Conditions that impair absorption in the distal small bowel and colon, such as coeliac disease, Crohn's disease, short bowel syndrome, or chronic diarrhoea.
    \item \textbf{Renal wasting:} Increased urinary excretion caused by osmotic diuresis in uncontrolled diabetes mellitus, hyperaldosteronism, or intrinsic renal tubular defects (such as Gitelman or Bartter syndromes).
    \item \textbf{Medication-induced loss:} Long-term use of proton pump inhibitors (PPIs, e.g. omeprazole), which inhibit active intestinal magnesium transport, or diuretics (e.g. frusemide, hydrochlorothiazide), which increase renal excretion.
    \item \textbf{Chronic alcohol misuse:} Alcohol acts as a direct renal toxin that impairs the tubular reabsorption of magnesium, compounded by poor nutritional intake and frequent gastrointestinal losses.
    \item \textbf{Intracellular redistribution:} Shifting of magnesium from extracellular to intracellular compartments, occurring in hungry bone syndrome following parathyroidectomy, or during the treatment of diabetic ketoacidosis with insulin.
\end{itemize}

\section*{Clinical Features}
The clinical features of hypomagnesaemia are primarily neuromuscular and cardiovascular, reflecting magnesium's role in stabilizing cell membranes and regulating electrical activity. Features include:
\begin{itemize}
    \item \textbf{Neuromuscular hyperexcitablity:} Muscle cramps, spasms, tremors, fasciculations (twitching), hyperreflexia, and in severe cases, painful tetany.
    \item \textbf{Positive latent tetany signs:} Chvostek's sign (ipsilateral facial muscle contraction upon tapping the facial nerve) and Trousseau's sign (carpopedal spasm induced by inflating a blood pressure cuff above systolic pressure), which are often secondary to concurrent hypocalcaemia.
    \item \textbf{Neurological disturbances:} Generalized weakness, lethargy, apathy, depression, irritability, and in severe cases, seizures, nystagmus, or delirium.
    \item \textbf{Cardiac arrhythmias:} Palpitations, premature ventricular contractions (PVCs), atrial fibrillation, and life-threatening ventricular arrhythmias.
    \item \textbf{Secondary electrolyte imbalances:} Refractory hypokalaemia and hypocalcaemia, presenting with their own respective symptoms, which cannot be corrected until the underlying magnesium deficiency is resolved.
    \item \textbf{Gastrointestinal features:} Anorexia, nausea, and vomiting, which can further exacerbate the deficiency.
\end{itemize}

\section*{Differential Diagnosis}
Because the features of magnesium deficiency overlap significantly with other mineral and metabolic disorders, differential diagnoses must be excluded:
\begin{itemize}
    \item \textbf{Hypocalcaemia:} Presents with identical signs of neuromuscular irritability (tetany, positive Chvostek's and Trousseau's signs), but is differentiated by low ionized calcium levels on biochemical testing (though the two conditions frequently coexist).
    \item \textbf{Hypokalaemia:} Presents with muscle weakness, cramps, and arrhythmias, distinguished by low serum potassium levels.
    \item \textbf{Vitamin D deficiency:} Leads to muscle weakness and hypocalcaemia, but is characterised by low 25-hydroxyvitamin D and elevated parathyroid hormone (PTH) levels.
    \item \textbf{Hyperventilation and respiratory alkalosis:} Causes transient carpopedal spasm and circumoral paresthesias due to a temporary decrease in ionized calcium, but is accompanied by hypocapnia and elevated arterial pH, with normal serum magnesium.
    \item \textbf{Primary neuromuscular disorders:} Conditions such as myasthenia gravis or peripheral neuropathies, which cause weakness or fasciculations but present with normal electrolyte panels.
\end{itemize}

\section*{When to Seek Medical Care}
Patients experiencing persistent muscle spasms, unexplained weakness, or numbness should consult a general practitioner for blood tests. Individuals taking long-term PPIs or diuretics should have their electrolytes monitored regularly.

\textbf{Call 000 (or local emergency services) for:}
\begin{itemize}
    \item Sudden onset of seizures or convulsions.
    \item Rapid, irregular, or pounding heartbeat (palpitations) associated with chest pain or dizziness.
    \item Severe difficulty breathing or swallowing (indicating potential laryngeal spasm).
    \item Loss of consciousness or sudden, severe confusion.
\end{itemize}

\section*{Diagnosis and Investigations}
Confirming magnesium deficiency involves biochemical assays, electrocardiography, and assessments of renal function:
\begin{itemize}
    \item \textbf{Serum magnesium level:} The primary diagnostic test, with a level $<0.7$ mmol/L confirming hypomagnesaemia. Severe deficiency is defined as a level $<0.4$ mmol/L.
    \item \textbf{Basic biochemistry panel:} Assays for serum potassium, calcium, phosphate, and creatinine to detect concurrent electrolyte deficiencies and evaluate renal function.
    \item \textbf{Electrocardiogram (ECG):} Essential for identifying cardiac conduction abnormalities, which typically include prolongation of the PR and QT intervals, widening of the QRS complex, flattening or inversion of T-waves, and U-waves.
    \item \textbf{24-hour urinary magnesium excretion:} Performed to differentiate renal wasting (urinary magnesium $>1.0$ mmol/day or fractional excretion $>2\%$) from gastrointestinal malabsorption or dietary insufficiency.
    \item \textbf{Intravenous magnesium loading test:} In patients with suspected tissue depletion despite normal serum levels, urinary excretion of less than 60--70\% of an infused magnesium load over 24 hours confirms a subclinical deficiency.
\end{itemize}

\section*{Management}
Treatment is determined by the severity of the deficiency, the presence of symptoms, and the patient's renal function:
\begin{itemize}
    \item \textbf{Oral magnesium supplementation:} For mild, asymptomatic deficiency, utilizing organic magnesium salts (such as magnesium citrate, aspartate, or lactate) which are better absorbed and cause less diarrhoea than inorganic salts like magnesium oxide.
    \item \textbf{Intravenous magnesium replacement:} For severe or symptomatic hypomagnesaemia ($<0.4$ mmol/L), or in the presence of seizures or arrhythmias, administered as intravenous magnesium sulfate (usually 8--16 mmol infused over several hours).
    \item \textbf{Emergency administration:} In the case of torsades de pointes, 8 mmol of magnesium sulfate is given as an immediate intravenous bolus over 1--2 minutes.
    \item \textbf{Correction of secondary cytopenias:} Intravenous or oral potassium and calcium replacement, which will only succeed once magnesium levels are restored.
    \item \textbf{Treatment of the cause:} Discontinuing or reducing offending drugs (such as PPIs or loop diuretics) where clinically safe, and treating underlying chronic bowel or metabolic conditions.
\end{itemize}

\section*{Complications}
Untreated magnesium deficiency can lead to severe cardiovascular and neurological outcomes:
\begin{itemize}
    \item \textbf{Torsades de pointes:} A life-threatening polymorphic ventricular tachycardia that can degenerate into ventricular fibrillation and sudden cardiac death.
    \item \textbf{Seizures and status epilepticus:} Triggered by uncontrolled neuronal hyperexcitability.
    \item \textbf{Laryngeal spasm:} Spasm of the vocal cords causing acute upper airway obstruction, which is a medical emergency.
    \item \textbf{Chronic hypocalcaemia and hypokalaemia:} Causing persistent neuromuscular and cardiac symptoms due to impaired PTH secretion and renal potassium wasting.
    \item \textbf{Skeletal fragility:} Long-term magnesium depletion impairs bone mineralisation and promotes osteoporosis.
\end{itemize}

\section*{Prognosis and Outcomes}
The prognosis for magnesium deficiency is excellent when identified and treated promptly. Serum levels typically normalise within a few days of starting replacement therapy, and neuromuscular irritability resolves rapidly. For patients with permanent risk factors (such as chronic malabsorption or an absolute requirement for PPI therapy), long-term, low-dose oral supplementation is highly effective at maintaining normal magnesium homeostasis. The prognosis is only poor if severe deficiency is unrecognized, leading to fatal cardiac arrhythmias.

\section*{Prevention and Self-Care}
Practical dietary and lifestyle measures to maintain healthy magnesium levels include:
\begin{itemize}
    \item \textbf{Dietary optimization:} Incorporate magnesium-rich foods into daily meals, including spinach, chard, pumpkin seeds, almonds, cashews, black beans, whole grains, avocados, and dark chocolate.
    \item \textbf{Moderate alcohol intake:} Limit alcohol consumption to prevent transient renal magnesium wasting and general nutritional depletion.
    \item \textbf{Supplementation awareness:} Take oral magnesium supplements only under medical supervision, especially in patients with pre-existing kidney disease, as excess magnesium can lead to toxic hypermagnesaemia.
    \item \textbf{Diabetes control:} Manage blood glucose levels carefully in type 2 diabetes to reduce osmotic diuresis and subsequent renal magnesium loss.
    \item \textbf{Medication monitoring:} Patients on long-term PPIs or diuretics should request annual electrolyte screening from their general practitioner.
\end{itemize}

\section*{Sources and Further Reading}
The following reputable organisations provide further information on magnesium deficiency:
\begin{itemize}
    \item Healthdirect Australia. \textit{Information on magnesium deficiency, symptoms, and prevention.} https://www.healthdirect.gov.au/
    \item Better Health Channel. \textit{Dietary minerals, including magnesium roles and sources.} https://www.betterhealth.vic.gov.au/
    \item NHS. \textit{Overview of essential minerals and hypomagnesaemia.} https://www.nhs.uk/
    \item Mayo Clinic. \textit{Electrolyte imbalances: symptoms and clinical management.} https://www.mayoclinic.org/
    \item MedlinePlus. \textit{Magnesium in the human diet and deficiency states.} https://medlineplus.gov/
\end{itemize}
